\section{Scope, conclusions, and verification status}

The source's public conclusions concern the Riemann zeta function and
\emph{Dirichlet} \(L\)-functions. Its entry point proves the five statements in
\src{Main}{\code{Main.lean}}, using the Mathlib-facing definitions in
\src{Statements}{\code{Statements.lean}}. This review uses GRH with that
Dirichlet scope; it does not extend the source's conclusion to all
automorphic or Artin \(L\)-functions.

The pinned readable tree contains 803 Lean files and 477,020 physical lines in
\code{src}. Physical lines include comments and blanks. A syntactic import
walk from \code{Main} reaches all 803 files. This does not show that every
declaration is needed: imports and proof-term dependencies are different
graphs. It does show that a claim of large savings from merely removing
unreachable imported files would be unsupported by this snapshot.
The archive contains the measurements and a reproducible audit script.

The supplied \guide{} explains the arithmetic determinant comparison and is
valuable for locating the conceptual bottleneck. The proposed changes below
were checked against the corresponding Lean definitions and hypotheses, not
only against the guide's prose.

\begin{center}
\begin{tabular}{@{}p{.34\textwidth}p{.60\textwidth}@{}}
\toprule
Component & Status in this delivery\\
\midrule
Exact power homogeneity of the common-multiplier height &
Kernel-checked standalone generalization of the actual definition.\\
Determinant recurrence; additive costs; triangular count &
Kernel-checked abstract lemmas. They assume the stated recurrence.\\
Specialization to \code{New.HeightLe} &
Insertion-ready candidate, explicitly not compiled in the full source tree.\\
Local exchange-to-cost conversion and full determinant migration &
Mathematical derivation and integration plan; not a completed patch.\\
One-vector branching norm bound &
Mathematical proof with existing source prerequisites identified;
source specialization not kernel-checked here.\\
The full RH/GRH project &
Read and analyzed at relevant interfaces; not rebuilt or independently
replayed in this session.\\
\bottomrule
\end{tabular}
\end{center}

The repository's \href{https://github.com/VladimirReshetnikov/ProveIt/blob/9050b04000647d588d0c78d945180abcd3b0bcf7/Analysis/GeneralizedRiemannHypothesis/Lean/grh/README.md}{README}
reports full builds and independent checker acceptance. Those are the
repository's reports, not new verification claims made here. For the new
standalone lemmas, compilation succeeded and the printed axiom sets contain
only \code{propext}, \code{Classical.choice}, and \code{Quot.sound}.
No revised proof of the five final theorems is being claimed.

The practical conclusion is two concrete, compatible simplifications of
substantial local arguments. I do not claim a measured percentage reduction
of the entire 477,020-line development. A net line count and build-time claim
require completing the integrations.

\section{The part of the proof that must survive}

\subsection{The arithmetic comparison}

Fix the progression modulus \(Q\), an admissible auxiliary prime \(d\), and the
quadratic character \(\chi_d\). The proof constructs a long chain of theta lifts
whose length satisfies \(m\asymp\sqrt{x}\). Its upper estimate makes a large
prime discrepancy produce many short vectors. Its spanning estimate ensures
that sufficiently many of those vectors are independent. The lower estimate
uses integrality, twist-dependent height bounds, and the product formula to
limit how short their combined volume can be.

Write \(d_T\) for the dimension of the target forms, to avoid confusing it
with the auxiliary prime. In the notation of the source, put
\[
 L=fa(\log b)^4,\qquad
 H=m k\,\operatorname{directionAmount}=\operatorname{mainAmount}.
\]
Theorem C is the statement that, for an independent family with
\(n\ge(1-4/b^2)d_T\), logarithmic squared-covolume allowance \(V_0\), and
index allowance \(\delta\),
\begin{equation}
 V_0+2\delta\ \ge\ d_T(H-3L).
 \label{eq:theoremC}
\end{equation}
This is
\src{New/Determinant/Twist/LongForms}{\code{integral_forms_long}}.
The coefficient of \(H\) matters. Losing a copy of the ordinary-direction
contribution changes the cancellation with the short-form theorem.

The three inputs are assembled in
\src{New/Reduction/MainInequality}{\code{MainInequality.lean}}.
The eventual one-sided prime-remainder estimate is then converted to the
Dirichlet zero-free half-plane in
\src{New/Reduction/OneSidedBound}{\code{OneSidedBound.lean}} and the
\src{New/Reduction/Zeros/Criterion}{zero-detection criterion}.
The half-plane statement and the functional equation produce the
public critical-line statements.

The main proposal changes how the determinant lower bound is organized.
It preserves \eqref{eq:theoremC}, so it can coexist with the stronger
spanning, short-form, torsion-selection, and analytic proposals already
in the repository.

\subsection{Why a chain cannot simply be treated as independent single lifts}

The source uses a triangular diagram of lifts and descents. Moving a descent
past a lift introduces a scalar. For a chain of length \(m\), the crossings are
indexed by
\[
 \{(i,j):1\le i<j\le m\}.
\]
The first arrows in the rows have the stable ranks needed by the Rallis
formula, whereas consecutive arrows of the long chain do not all have that
property. Thus a valid simplification retains the square laws, their adjoint
laws, and the induction proving injectivity.

The relevant interfaces are
\src{New/Determinant/Chain/TriangleNorms}{\code{TriangleNorms}},
\src{New/Determinant/Chain/ChainNorms}{\code{ChainNorms}}, and
\src{New/Determinant/Chain/NormOfChain}{\code{NormOfChain}}.
The simplification proposed here removes unnecessary information carried
\emph{through} that induction; it does not replace the induction by an
unjustified multiplication of single-lift volumes.
